% =========================================================
% Proof: Q is Dense in R
% Source: notes/metric-spaces/notes-dense-separable.tex
% Difficulty: Easy
% =========================================================

\subsection*{$\mathbb{Q}$ is Dense in $\mathbb{R}$}
\label{prf:Q-dense-R}

\begin{remark*}[Return]
$\leftarrow$ \hyperref[prop:Q-dense-R]{Back to Proposition in Notes}
\end{remark*}

\begin{proof}
\Claim $\mathbb{Q}$ is dense in $(\mathbb{R}, |\cdot|)$: for all $x \in \mathbb{R}$
and all $r > 0$, the ball $B_r(x) = (x-r, x+r)$ contains a rational.

\Given $x \in \mathbb{R}$ and $r > 0$.

\Goal To find $q \in \mathbb{Q}$ with $|q - x| < r$.

\Strategy This is exactly the density of $\mathbb{Q}$ in $\mathbb{R}$
from the real analysis foundations, proved via the Archimedean property.
State which result from foundations you are using, then apply it directly.

% -------------------------------------------------------
% YOUR PROOF HERE
% -------------------------------------------------------

\end{proof}

\begin{remark*}[Proof shape]
Direct application of the Archimedean property: between any two real
numbers there is a rational.
\end{remark*}

\begin{remark*}[Dependencies]
\begin{itemize}
  \item Archimedean property: for any $x \in \mathbb{R}$ and $r > 0$,
        there exists $q \in \mathbb{Q}$ with $|q - x| < r$.
  \item Proposition~\ref{prop:dense-equiv}(ii): density via open balls.
\end{itemize}
\end{remark*}
